By \(\text{thm:rank-adaptive-estimation}\),
\[
 \sqrt n\lVert\widehat A_n-A\rVert_{\max}=O_{P_\theta}(1)
\tag{1}
\]
uniformly on \(\mathcal M_{\Delta,p}(\kappa)\), since the maximum norm is bounded by the Frobenius norm.  We convert \(\text{ass:edge-gap}\), which is stated for the signed jump, to a uniform amplitude margin.  Acyclicity gives \(A^p=0\) and therefore
\[
 \lambda=(I_p-A/\beta)^{-1}\mu
 =\sum_{d=0}^{p-1}(A/\beta)^d\mu.
\tag{2}
\]
On the margin class, \(\beta\geq\underline\beta=-\Delta^{-1}\log(1-\kappa)\), \(\lVert A\rVert_{\mathrm{op}}\leq\kappa^{-1}\), and \(\lVert\mu\rVert_2\leq\kappa^{-1}\).  Hence
\[
 \max_j\lambda_j\leq\lVert\lambda\rVert_2
 \leq\Lambda_\kappa:=\kappa^{-1}
 \sum_{d=0}^{p-1}(\kappa^{-1}/\underline\beta)^d<\infty.
\tag{3}
\]
For a true edge \(j\to i\), acyclicity excludes \(i\to j\), and the skew-jump identity in \(\text{thm:two-lag-exact-recovery}\) gives \(Z_{ij}=A_{ij}\lambda_j\).  Thus \(\text{ass:edge-gap}\) and (3) imply
\[
 A_{ij}=|Z_{ij}|/\lambda_j\geq\kappa/\Lambda_\kappa.
\tag{4}
\]
For an absent directed edge, \(A_{ij}=0\).  For all sufficiently large \(n\), \(a_n<\kappa/(2\Lambda_\kappa)\).  On the event
\[
 \lVert\widehat A_n-A\rVert_{\max}
 <\min\{a_n,\kappa/(2\Lambda_\kappa)\},
\tag{5}
\]
there are neither false inclusions nor false exclusions, so \(\widehat G_n=G(A)\).  By (1) and the stated condition \(\sqrt n\,a_n\to\infty\), the probability of the complement of (5) tends to zero uniformly.  This proves the proposition.